\chapter{总结与展望}

\section{全文工作总结}
本文围绕“基于视觉语言预训练与提示学习的无监督医学图像检测、分割及数字病理分析研究”这一主题，系统开展了四项相互关联的研究工作，并形成了一条从基础视觉识别问题逐步走向临床病理分析问题的连续主线。

首先，本文从图像级弱监督细胞核实例分割问题出发，提出循环学习框架，解决在监督极其有限条件下如何构造有效实例学习信号的问题。这一工作奠定了全文“监督补偿”的基础逻辑。其次，本文将研究扩展到视觉语言预训练模型，提出自动属性提示驱动的零样本细胞核检测方法，说明即使缺乏专门检测标注，模型仍可借助跨模态先验获得起始检测能力。再次，本文针对文本提示表达不足的问题，提出视觉文本化方法，使图像示例能够以提示形式进入视觉语言检测器，实现从文本提示向图像提示的扩展。最后，本文将前述方法能力迁移到肺鳞癌数字病理分析中，尝试构建联合残余肿瘤负荷和肿瘤周围免疫细胞密度的临床量化分析框架。

从技术脉络来看，本文的四项工作构成一条连续的研究链条。第二章解决的是监督不足下的学习信号构造问题；第三章进一步把外部视觉语言先验引入医学检测；第四章则继续讨论当纯文本提示不足时如何利用图像示例增强模型；第五章最终说明这些方法研究为什么值得做，即它们能够服务于更高层次的临床问题分析。正是这种问题链的连续性，使整篇学位论文的叙事结构更加统一。

\section{主要贡献}
综合全文，本文的主要贡献可以归纳为以下几点。

第一，提出了适用于细胞核实例分割的循环学习框架，说明在仅有图像级监督的条件下，通过前后端闭环优化可以逐步逼近实例级任务目标。第二，提出了自动属性提示与自训练联合的零样本细胞核检测框架，为无专门检测标注的医学视觉任务提供了新的解决路径。第三，提出了视觉文本化方法，将图像提示映射到语言空间，从而增强模型在开放词汇和跨域条件下的适应能力。第四，构建了面向肺鳞癌疗效评估的数字病理分析框架，把前端识别结果组织成具有临床解释意义的病理量化指标。

\section{研究意义}
从理论意义上看，本文尝试将弱监督学习、视觉语言预训练、提示学习和数字病理分析串联起来，说明这些看似分散的研究方向可以围绕“监督不足与泛化不足”这一核心问题形成统一的方法体系。本文既关注模型性能，也关注学习信号来源、提示机制构造和结果解释性，因此在方法论上具有一定连续性。

从实践意义上看，本文的研究为低标注成本医学图像分析提供了若干可行路径。无论是通过弱监督和自训练减少精细标注依赖，还是通过视觉语言模型和图像提示提升跨域适配能力，最终目标都是提高方法在真实医学场景中的可扩展性。特别是肺癌数字病理分析工作进一步说明，前端识别能力可以转化为更高层次的临床量化价值。

\section{局限性与不足}
尽管本文取得了若干阶段性成果，但仍存在明显不足。首先，视觉语言预训练模型在医学场景中的迁移能力仍受到预训练语料分布限制，自动属性提示和图像提示虽然可以缓解问题，但尚未从根本上解决医学语义鸿沟。其次，图像提示分支之间的平衡和稳定性仍然较为敏感，提示学习方法在不同数据集上的表现可能受支持样本选择影响较大。再次，在临床应用层面，数字病理分析工作虽然初步展示了量化指标价值，但仍需更大规模、多中心验证以进一步支撑临床结论。

\section{研究展望}
未来工作可以从以下几个方向进一步推进。

第一，研究更自动化、更稳健的提示生成机制，使文本提示、图像提示和结构提示能够根据任务特点动态协同。第二，探索更强的域泛化和领域自适应策略，以增强模型在多中心、多染色和复杂真实场景中的稳定性。第三，尝试将检测、分割、分类和病理量化统一到多任务框架中，减少分阶段误差积累。第四，在临床方向上，继续推动算法结果与病理知识、结局分析和多模态临床数据融合，从而构建更加完整的 AI 辅助病理分析系统。

\section{结语}
总体而言，本文尝试构建一条从弱监督识别、零样本检测、图像提示学习到临床数字病理分析的连续研究路线。尽管仍有许多工作需要进一步完善，但现有成果已经说明：视觉语言预训练与提示学习为医学图像分析提供了新的研究范式，也为降低标注依赖、提升跨域泛化和增强临床解释性打开了新的空间。
